
\begin{abstract} 
	
	The ability to characterize the functional behavior of an individual protein in a variety of different contexts is one of the cornerstones of carrying out biological investigations. Its importance is especially apparent in the context of understanding life at the molecular level, investigating disease, the development of drugs to cure disease, and the manner in which evolution is modulated by changes in genes. 
	
Although a growing number of tools for the task of interrogating a protein's function are at the disposal of experimental biologists, characterizing the exponentially expanding set of known sequences requires the application of \textit{in silico} techniques.
 
In this dissertation novel methods for the prediction of protein function from protein sequence and structure data are detailed. Current techniques for the evaluation of function prediction are reviewed and several drawbacks are addressed through the introduction of new information content-based metrics that attempt to provide comprehensive assessment of the performance of a given prediction method. Finally, the influences of speciation and gene duplication on the evolution of protein function is investigated by comparing the functional similarity of different classes of homologous sequences.
\end{abstract}
